\documentclass[conference]{IEEEtran}
\IEEEoverridecommandlockouts

% Citation management
\usepackage{cite}

% Text and list formatting
\usepackage{enumitem}
\usepackage{amsmath,amssymb,amsfonts}

% Algorithms (optional but common in empirical papers)
\usepackage{algorithm}
\usepackage{algorithmic}

% Graphics and visual elements
\usepackage{graphicx}
\usepackage{textcomp}
\usepackage{xcolor}
\usepackage{url}

% Figure and table enhancements
\usepackage{booktabs}
\usepackage{threeparttable}

% Diagrams (optional)
\usepackage{tikz}
\usetikzlibrary{arrows,positioning,fit,calc}

\begin{document}

\title{[Empirical Paper Title]}

\author{%
\IEEEauthorblockN{[Author Name(s)]}
\IEEEauthorblockA{[Affiliation(s)] \\
Correspondence: [email@example.com]}
}

\maketitle

\begin{IEEEkeywords}
[keyword 1], [keyword 2], [keyword 3], [keyword 4]
\end{IEEEkeywords}

\begin{abstract}
[Abstract placeholder. State the problem, the core idea, the experimental setting,
the key outcome only if verified, and the main implication. Avoid citations here.]
\end{abstract}

\section{Introduction}
[State the problem, why existing work is insufficient, the main contribution, and
the paper roadmap. Keep every non-trivial claim citation-backed.]

\subsection{Contributions (outline)}
\begin{itemize}[leftmargin=*]
\item [Contribution 1]: TBD (map to claim IDs in `brief/contribution-map.yaml` if available).
\item [Contribution 2]: TBD.
\item [Contribution 3]: TBD.
\end{itemize}

\section{Related Work}
[Summarize the closest prior methods/settings and clearly distinguish this paper's
claim from them.]

\section{Method}
[Describe the proposed method, formulation, architecture/training choices, and
the rationale behind each design decision.]

\subsection{Problem Setup}
[Inputs, outputs, retrieval setting, and assumptions.]

\subsection{Approach Overview}
[High-level pipeline diagram placeholder.]

\begin{figure}[t]
\centering
\fbox{\parbox{0.95\columnwidth}{\vspace{18mm}\centering Pipeline diagram (placeholder)\vspace{18mm}}}
\caption{System overview (placeholder).}
\label{fig:pipeline}
\end{figure}

\section{Experimental Setup}
[Define datasets, baselines, metrics, hardware/runtime if relevant, and the exact
evaluation protocol.]

\subsection{Datasets}
[List datasets and splits (planned/verified).]

\subsection{Baselines}
[List baselines and versions.]

\subsection{Metrics}
[Define primary and secondary metrics.]

\section{Results}
[Main results, ablations, and analysis go here. Use explicit placeholder text
instead of invented numbers when results are not yet verified.]

\subsection{Main Results}
[Verified or placeholder-safe comparison to baselines.]

\begin{table}[t]
\centering
\caption{Main results (placeholder). Use ``--'' for unverified values.}
\label{tab:main_results}
\begin{tabular}{lccc}
\toprule
Method & Metric 1 & Metric 2 & Cost \\
\midrule
Baseline A & -- & -- & -- \\
Baseline B & -- & -- & -- \\
Ours & -- & -- & -- \\
\bottomrule
\end{tabular}
\end{table}

\subsection{Ablation Studies}
[What is removed/changed and why it matters.]

\subsection{Error Analysis / Robustness}
[Failure cases, robustness checks, or uncertainty analysis.]

\section{Limitations}
[Explicit limitations, open risks, and evidence not yet collected.]

\section{Conclusion}
[Summarize only claims that are already supported by the draft's evidence state.]

\label{ReferencesStart}
\bibliographystyle{ieeetr}
\bibliography{ref}

\end{document}
